\chapter{Topic: Exception Safety Guarantees}

\begin{aside}
This chapter continues Lecture 21 and covers material from the remainder of the
term (not in the original lecture notes).
\end{aside}

\section{Levels of Exception Safety}
Lecture 21 made sure resources are released when an exception passes through. But
what state is an object left in? There are three standard levels of guarantee a
function can offer:

\begin{definition}
\begin{enumerate}
  \item \textbf{Basic guarantee}: if an exception occurs, the program is in
        \emph{some} valid but unspecified state. Nothing is leaked, no
        invariants are broken.
  \item \textbf{Strong guarantee}: if an exception is raised while executing
        \code{f}, the state of the program is as if \code{f} had not been
        called (``commit or roll back'').
  \item \textbf{No-throw guarantee}: \code{f} never throws an exception, and
        always accomplishes its task.
\end{enumerate}
\end{definition}
A function that gives none of these is not exception safe at all.

\section{An Example}
\begin{lstlisting}[style=cpp]
class A { ... };
class B { ... };
class C {
    A a;
    B b;
  public:
    void f() {
        a.g();   // may throw (strong guarantee)
        b.h();   // may throw (strong guarantee)
    }
};
\end{lstlisting}
Is \code{C::f} exception safe?
\begin{itemize}
  \item If \code{a.g()} throws: nothing has happened yet (\code{g} has the
        strong guarantee). OK.
  \item If \code{b.h()} throws: the effects of \code{h} are undone, but
        \code{a.g()} has already happened and is not undone. So \code{f} would
        need to undo \code{g}'s effects --- very hard or impossible if \code{g}
        had non-local side effects (printing, writing a file).
\end{itemize}
So \code{f} offers only the basic guarantee (at best).

\subsection{Achieving the Strong Guarantee}
If \code{A::g} and \code{B::h} have no non-local side effects, we can do the
work on \emph{copies} and then swap them in:
\begin{lstlisting}[style=cpp]
class C {
    A a;
    B b;
  public:
    void f() {
        A atemp = a;    // may throw: fine, nothing changed yet
        B btemp = b;
        atemp.g();      // may throw: fine, only the copies changed
        btemp.h();
        a = atemp;      // what if copy assignment throws?
        b = btemp;
    }
};
\end{lstlisting}
The last two lines are the problem: if \code{b = btemp} throws, \code{a} has
already been modified. We need the ``commit'' step to be no-throw. Assigning
\emph{pointers} can't throw --- so use the pImpl idiom:
\begin{lstlisting}[style=cpp]
struct CImpl {
    A a;
    B b;
};

class C {
    std::unique_ptr<CImpl> pImpl;
  public:
    void f() {
        auto temp = std::make_unique<CImpl>(*pImpl);  // copy: may throw, fine
        temp->a.g();                                  // may throw, fine
        temp->b.h();                                  // may throw, fine
        std::swap(pImpl, temp);                       // no-throw commit
    }   // temp (the old state) is destroyed here
};
\end{lstlisting}
Now \code{f} has the strong guarantee (assuming \code{g} and \code{h} have no
non-local side effects). This is the same idea as copy-and-swap for
\code{operator=}.
\begin{center}
\begin{tikzpicture}[>=Stealth,font=\small,c/.style={draw,fill=blue!5,minimum width=1.4cm,minimum height=0.65cm}]
  \node[c] (pi) at (0,0) {\code{pImpl}};
  \node[c] (t) at (0,-1.2) {\code{temp}};
  \node[draw,fill=gray!12,minimum width=2.2cm,minimum height=0.7cm] (old) at (3.6,0) {old \code{a}, \code{b}};
  \node[draw,fill=green!12,minimum width=2.2cm,minimum height=0.7cm] (new) at (3.6,-1.2) {modified copy};
  \draw[->] (pi) -- (old); \draw[->] (t) -- (new);
  \draw[<->,thick,blue!60!black] (-0.9,0) to[bend right=40] node[left,font=\scriptsize]{\code{swap}} (-0.9,-1.2);
  \node[font=\scriptsize,align=left,anchor=west] at (5.0,-0.6) {all throwing work happens on the copy;\\ the swap of two pointers cannot throw};
\end{tikzpicture}
\end{center}

\begin{important}[Rule of thumb]
Do everything that might throw on the side (on copies/temporaries), then commit
with operations that cannot throw (pointer assignment, swap of built-in types,
moves marked \code{noexcept}).
\end{important}

\section{Exception Safety and \code{vector}}
\code{vector} manages a heap array with RAII: when a vector goes out of scope,
its dtor frees the array and calls the dtor of every element.

\begin{itemize}
  \item \code{vector<C>}: owns the \code{C} objects; they are destroyed with
        the vector.
  \item \code{vector<C*>}: the dtor frees the array of pointers, but does
        \emph{not} delete the objects pointed to (pointers have no dtor). If the
        vector should own them, you must delete them yourself:
\begin{lstlisting}[style=cpp]
for (auto p : v) delete p;
\end{lstlisting}
  \item \code{vector<unique\_ptr<C>>}: owns the objects, and frees them
        automatically, even during unwinding. No explicit cleanup loop needed.
\end{itemize}

\subsection{\code{emplace\_back} and \code{noexcept}}
\code{vector::emplace\_back} (and \code{push\_back}) offer the strong guarantee:
if it throws, the vector is unchanged. When the array is full, it must
\begin{enumerate}
  \item allocate a bigger array,
  \item copy the elements from the old array to the new one,
  \item delete the old array.
\end{enumerate}
If a copy throws midway (step 2), the new array is simply discarded and the old
one is intact: strong guarantee.

But copying is expensive; moving would be much faster. However, if a \emph{move}
throws midway, some elements of the old array have already been gutted, and the
original state can't be restored. So \code{vector} only moves elements if the
move ctor is guaranteed not to throw; otherwise it copies.
\begin{center}
\begin{tikzpicture}[>=Stealth,font=\small,x=0.8cm,c/.style={draw,minimum width=0.8cm,minimum height=0.6cm}]
  \node[anchor=east,font=\scriptsize] at (-0.6,0) {old array (full)};
  \foreach \k in {0,...,3} \node[c,fill=gray!12] (o\k) at (\k,0) {$e_\k$};
  \node[anchor=east,font=\scriptsize] at (-0.6,-1.4) {new array};
  \foreach \k in {0,...,3} \node[c,fill=green!10] (n\k) at (\k,-1.4) {$e_\k$};
  \foreach \k in {4,...,7} \node[c] at (\k,-1.4) {};
  \foreach \k in {0,...,3} \draw[->] (o\k) -- (n\k);
  \node[font=\scriptsize,anchor=west,align=left] at (4.4,0) {copied, or moved if the move ctor\\ is \code{noexcept}; then the old array is freed};
\end{tikzpicture}
\end{center}

\begin{important}[Advice]
Declare move constructors and move assignment operators (and \code{swap}) as
\code{noexcept} when they cannot throw --- which is almost always.
\begin{lstlisting}[style=cpp]
class C {
  public:
    C(C &&other) noexcept;
    C &operator=(C &&other) noexcept;
};
\end{lstlisting}
Otherwise \code{vector} (and other library code) will silently fall back to the
slower copy. Destructors are implicitly \code{noexcept}.
\end{important}

\section{Exception-Safe \code{operator=}}
Compare the versions of \code{Node::operator=} from Lectures 10--11:
\begin{center}
\begin{tabular}{@{}lp{0.55\textwidth}@{}}
\toprule
Version & Guarantee \\
\midrule
\code{delete next;} then \code{new} & none: leaves a dangling \code{next} if
\code{new} throws \\
\code{new} into temp, then \code{delete} old & strong \\
copy-and-swap & strong (if the copy ctor is safe and \code{swap} is no-throw) \\
\bottomrule
\end{tabular}
\end{center}

\section*{Summary}
\begin{itemize}
  \item Basic: valid state, no leaks. Strong: all-or-nothing. No-throw: never
        fails.
  \item A sequence of strong operations is not automatically strong.
  \item Achieve strong guarantee by working on copies and committing with
        no-throw operations (pImpl + swap).
  \item Use \code{vector<unique\_ptr<C>>} for owned polymorphic objects.
  \item Mark moves and swaps \code{noexcept} so containers can use them.
\end{itemize}
